\cotefile{130}
\scanpage{37}
si $K'/K$ est radicielle, et \deleted{\unsure{que est}}\inserted{plus gén.} \deleted{telle que} si $K'/K$ est quasi-galoisienne de \unsure{groupe|grpe} $\pi$, on a
\[
\widehat{\Phi}(X/K)\simeq\widehat{\Phi}(X'_{K'}/K')^{\pi}.
\]

\uline{Fonctions et parties $M_K$-bornées.}

\deleted{Une partie $E$ de $X$}

Soit $E$ un ens muni d'une application $\pi:E\longrightarrow M_{\overline K}$. Une fonction $\varphi$ sur $E$ est dite $M_{\overline K}$-\uline{bornée} \inserted{(rel à $\pi$)} s'il existe un $M_K$-\uline{diviseur} $c$ \deleted{\unsure{de} et une} \inserted{tel que}
\[
|\varphi(\unsure{\xi|\zeta})|\leq c(\pi(\xi))
\]
i.e. s'il existe une partie finie $S$ de $M_K$ \deleted{et un \unsure{entier}} telle que
\[
\varphi(\xi)=0\qquad\text{si}\quad \pi(\xi)|K\notin S
\]
et si de plus $\varphi$ est \uline{bornée}. \deleted{\unsure{Soit}}

% slanted
\margin{\uline{$X$ loc. de t.f. sur $K$}}
Soit $E\inserted{{}=\bigcup_{v\in M_{\overline K}} E_v\times\{v\}}$ une partie de $X(\overline K)\times M_{\overline K}$, on dit que $E$ est $M_K$-bornée s'il \deleted{existe une} \inserted{est \unsure{contenu|contenue}} dans une réunion finie d'ens de la forme $E(U;x_1,\ldots,x_{\unsure{n|m}};c)$, où $U$ est un ouvert
% archive sheet number: 36
\scanpage{38}
affine de $X$, $\underline{x}=(x_1,-,x_{\unsure{n|m}})$ est un syst. de générateurs de son anneau affine sur $K$, $c$ un $M_K$-diviseur, et $E(\text{---})$ est l'ens des $(P,v)\in\deleted{X}U(\overline K)\times M_{\overline K}$ tels que
\[
v(\underline{x}(P))=\mathop{\mathrm{Inf}}\limits_{1\leq i\leq n} v(x_i(P))\geq c(v)
\]
$\bigl\{$ i.e. \inserted{essentiellement} il existe une partie finie $S$ de $M_K$ telle que $E$ est formé des $(P,v)$ \inserted{$\in U(\overline K)\times M_{\overline K}$} tels que
\[
\begin{aligned}
&\text{a) }P\text{ soit entier pour }\underline{x},v\quad\text{si}\quad v|K\in M_K-S\\
&\text{b) }v(\unsure{\underline{x}(P)|x(P)})\geq-M\quad\text{si}\quad v|K\in S\quad\bigl\}.
\end{aligned}
\]

% slanted
\margin{\uline{N.B.} La notion de partie bornée est inv. par ext. finie du corps de base. La fonctorialité pour $X\to Y$ n'est pas claire, aussi il y aurait peut-être lieu de remplacer $U\unsure{\hookrightarrow|\subset}X$ par $U\to X$ ----}
% slanted
\margin{\uline{Question.} Dans le cas ``global'', une partie ``$M_K$-bornée'' de $\mathbb{P}^{r}_K$ est-elle \uline{de} \uline{hauteur} \uline{bornée} ?}

\uline{Fonctions admissibles}

Une fonction $\varphi:X(\overline K)\times M_{\overline K}\longrightarrow\mathbb R$ est dite \uline{admissible} si elle est \deleted{\unsure{bor}} $M_K$-bornée sur toute partie $M_K$-bornée. Comme les ens $\{P\}\times M_{\overline K}$, où $P\in X(\overline K)$, sont $M_K$-bornés, on voit que \deleted{si} pour tt $P\in\deleted{X}U(\overline K)_{\unsure{\theta|{}}}$ fixé, la fonction
\[
\widetilde{\varphi}(P)=(v\mapsto\varphi(P,v)):M_{\overline K}\longrightarrow\mathbb R
\]
a la vertu d'être \unsure{majorée|majoré} \inserted{en module} par une fonction de la forme $c(\unsure{w|v})$, avec
% archive sheet number: 37
\scanpage{39}
$c$ un $M_K$-diviseur positif. D'ailleurs, \inserted{Supposons $\varphi$ rat $/K$ i.e.\textsuperscript{\unsure{b|6}}} si $P$ est rationnel sur \inserted{l'ext. finie} $K'$ de $K$, alors $\widetilde{\varphi}(P)$ est rat $/K'$ i.e. se factorise en $M_{\overline K}\longrightarrow M_{K'}\longrightarrow\mathbb R$, où $M_{K'}\unsure{-|\to}\mathbb R$ est ici un $K'$-diviseur. Donc
% slanted
\margin{\textsuperscript{b} inv. par $\mathrm{Gal}\,\overline K/K=\Gamma$}
\[
\left[\ \text{posant }\operatorname{Div}(M_{\overline K})=
\underset{\substack{K'/K\\\text{sous-ext. finie}\\\text{de }\overline K/K}}{\unsure{\underline{\lim}|\varinjlim}}
\operatorname{Div}(M_{K'})\ \right]
\]
on trouve une application
\[
\widetilde{\varphi}:X(\overline K)\unsure{-|\longrightarrow}\operatorname{Div}(M_{\overline K})
\]
qui a tt au moins les vertus \unsure{suivant|suivantes}

a) $\widetilde{\varphi}$ est compatible avec les op de $\Gamma=\mathrm{Gal}\,\overline K/K$

b) $\widetilde{\varphi}$ transforme tte partie $M_K$-bornée $F$ de $X(\overline K)$ [i.e. une partie $F$ telle que $F\times M_{\overline K}$ soit $M_K$-borné] en une partie \uline{bornée} de $M_{\overline K}$-diviseurs [bornée au sens de la relation d'ordre \unsure{qu'on aura|on aura|au sens :} $\exists\,c\in\operatorname{Div}(M_K)$ avec $|\widetilde{\varphi}(P)|\leq c$ pour tt $P\in F$].
% archive sheet number: 38
\scanpage{40}
\uline{Exemples de} $\varphi$ \uline{admissibles}.

Soit $f\in\Gamma(X,\unsure{\mathcal O|\underline O}_X^{\times})$, posons
\[
\varphi_f(P,v)=v(f(P))\qquad\text{si}\quad P\in X(\overline K),\,v\in M_{\overline K}
\]
C'est une fonction \uline{invariante par} $\Gamma$, je dis qu'elle est \uline{admissible}. \deleted{Il suffit de montrer que si $X$ est affine, \deleted{\unsure{comme}} $K[X]$ engendré par $x_1\,-\,x_n$, \inserted{$c\in\operatorname{Div}(M_K)$,} alors $\varphi_f$ est $M_K$-bornée sur $E[X;x_1,-x_n,c]$ i.e. $\exists\,c'\in\operatorname{Div}(M_K)$ tel que}
\[
\deleted{\mathop{\mathrm{Inf}}\limits_i v(x_i(P))\geq c(v)\Longrightarrow |v(f(P)|\leq c'(\unsure{v|w}).}
\]

Par fonctorialité, on est ramené au cas où $X=\mathrm G_{m\,K}$, \deleted{et à vérifier} \inserted{$f$ la section canonique $x$,} \deleted{que} Dans ce cas, \deleted{pour tte} une partie $M_K\unsure{+|-}$bornée de $X(\overline K)\times M_{\overline K}$ \inserted{$=\overline K^*\times M_{\overline K}$} est contenue dans une partie de la forme $E(M,S)$, $S$ partie finie de $M_K$, \deleted{$M\unsure{\geq|>}0$,} $M$ un nb $>0$, formé des $(\unsure{x|P},v)$ tels que
\[
\begin{aligned}
&\text{a) }v(x)\ \deleted{\text{est entier pour }v}\inserted{{}=0\quad\text{i.e. }x\text{ entier}}\quad\text{si}\quad v|K\notin S\\
&\text{b) }\deleted{\unsure{x|v}}\,|v(x)|\leq M\quad\text{si}\quad v|K\in S.
\end{aligned}
\]
% archive sheet number: 39
\scanpage{41}
Or on a
\[
\varphi_f(x,v)=v(x)\qquad\deleted{\text{qui est}}
\]
d'où aussitôt la conclusion.

\uline{Faisceautisation} Pour tt ouvert $U$ de $X$, soit $\underline{\Phi}_X(U)$ l'ens des fonctions admissibles $\Gamma$-\unsure{invariantes} sur $U(\overline K)\times M_{\overline K}$.

Pour $U$ variable, on trouve ainsi un \uline{préfaisceau} \inserted{$\underline{\Phi}_X$} sur $X$, je dis que c'est un \uline{faisceau}. Évidemment, c'est un ``préfaisceau séparé'', reste à montrer que si une fonction
\[
\varphi:X(\overline K)\times M_{\overline K}\unsure{-|\longrightarrow}\mathbb R
\]
invariante par $\Gamma$ est telle que l'on a $\varphi|U_i(\overline K)\times M_{\overline K}$ admissible pour tt $i$ (où $(U_i)$ est un rec. ouvert de $X$) alors $\varphi$ est admissible. On est ramené au cas où $X$ est affine, et il semble qu'on ait besoin de savoir que \inserted{\unsure{déjà|alors|la réunion}} \uline{tte partie} $M_K$-\uline{bornée} \uline{de} $X(\overline K)\times M_{\overline K}$
% archive sheet number: 40
\scanpage{42}
% slanted
\margin{\uline{à vérifier}}
\uline{est contenue dans une réunion d'images} \uline{de parties} $M_K$-\uline{bornées de} $U_i(\overline K)\times M_{\overline K}$.

Admettons ce point. On aura, par l'\uline{exemple précédent}, un hom de faisceaux
\diagram{\begin{tikzcd}[ampersand replacement=\&]
\underline{O}_X^*\arrow[r,"{\substack{f\mapsto\varphi_f\\ \deleted{\unsure{\xi|f}}}}"] \& \underline{\Phi}_X
\end{tikzcd}}
d'autre part on a un hom de faisceaux
\diagram{\begin{tikzcd}[ampersand replacement=\&]
{\deleted{R}\unsure{O|\underline O|\underline{\mathcal O}}}_X^*\arrow[r,hook,"{\substack{\unsure{f\leftarrow\!\vert f|f\mapsto f}\\ \deleted{\xi}}}"] \& \underline{R}_X^{\times}
\end{tikzcd}}
où $R_X^*$ est le faisceau des fonctions rationnelles sur $X$ [N.B. nous supposons $X$ réduit ou $=$ séparable $/K$, pour fixer les idées]. Nous voulons ``une généralisation commune des fonctions admissibles et des fonctions rationnelles'', et définissons pour ceci le \uline{faisceau des quasi-fonctions} comme le faisceau
\[
\begin{aligned}
\underline{\mathrm{QF}}_X&=\underline{\Phi}_X\mathop{\tikz[baseline=(op.base)]{\node[inner sep=0pt] (op) {$\unsure{\amalg|\times}$};\draw[overlay,thin] (-.33,-.22) .. controls (.09,-.33) and (-.04,.48) .. (.35,.49);}}\limits_{O_X^*}\underline R_X^{\times}
&&\text{(somme amalgamée)}\\
&\simeq\underline{\Phi}_X\mathbin{\tikz[baseline=(op.base)]{\node[inner sep=0pt] (op) {$\times$};\draw[overlay,thin] (-.33,-.22) .. controls (.09,-.33) and (-.04,.48) .. (.35,.49);}}\underline R_X^*/\unsure{\Theta|\theta}(\underline O_X^*)
\end{aligned}
\]
% archive sheet number: 41
\scanpage{43}
où
\[
\unsure{\Theta|\theta}:\underline O_X^*\longrightarrow\underline{\Phi}_X\mathbin{\tikz[baseline=(op.base)]{\node[inner sep=0pt] (op) {$\times$};\draw[overlay,thin] (-.33,-.22) .. controls (.09,-.33) and (-.04,.48) .. (.35,.49);}}R_X^*
\]
est défini par
\[
\Theta(f)=(-\varphi_f\mathpunct{\tikz[baseline=(op.base)]{\node[inner sep=0pt] (op) {$,$};\draw[overlay,thin] (-.33,-.22) .. controls (.09,-.33) and (-.04,.48) .. (.35,.49);}}f^{\deleted{\unsure{-1|1}}})
\]
On a donc deux suites exactes
\diagram{\begin{tikzcd}[ampersand replacement=\&,column sep=large]
\underline{\Phi}_X\arrow[r] \& \underline{\mathrm{QF}}_X\arrow[r,"{\mathrm{div}}"] \& \overset{\simeq\,\underline{R}_X^*/\underline{O}_X^*}{\underline{\mathrm{Div}}_X}\arrow[r]\arrow[rrd,dash,bend right=12] \& 0 \& \\ \& \& \& \& \text{diviseurs de \uline{Cartier}}
\end{tikzcd}}
\diagram{\begin{tikzcd}[ampersand replacement=\&,column sep=large]
\underline R_X^*\arrow[r,"{f\mapsto\unsure{\zeta_f|\xi_f|\{f\}}}"] \& \underline{\mathrm{QF}}_X\arrow[r] \& \underline{\Phi}_X/\operatorname{Im}\underline{O}_X^*\arrow[r] \& 0
\end{tikzcd}}

Les \inserted{sections sur $X$ du} noyau de $\underline{\Phi}_X\unsure{-|\to}\underline{\mathrm{QF}}_X$ est formé des fonctions admissibles $\varphi$ telles qu'il existe $f\in\Gamma(X,\underline O_X^*)$ avec $(\deleted{\unsure{0}}1,\varphi)=(f,-\varphi_f)$, ce qui exige $f=1$ et $-\varphi_f(\unsure{(P)v|(P)f|P,v})=v(f(P))=0$ d'où $\varphi=0$, donc $\underline{\Phi}_X-\underline{\mathrm{QF}}_X$ est injectif::

\uline{les quasi-fonctions de diviseur nul sont les fonctions admissibles sur} $X$.

Les \inserted{sections sur $X$ du} noyau de $\underline R_X^*-\underline{\mathrm{QF}}_X$ est formé des $f\in{\inserted{\Gamma R_X^*}}$ telles qu'il existe $g\in\Gamma(X,\underline O_X^{\unsure{\times|*}})$, t.q. $(f,0)=(g,-\varphi_g)$. On doit avoir
% archive sheet number: 42
\scanpage{44}
$\varphi_g=0$ i.e. $v(g(P))=0$ pour tt $P\in X(\overline K)$, $v\in M_{\overline K}$, donc $g(P)$ doit être une $M_{\overline K}$-unité pour tout $P\in X(\overline K)$.

Si $g$ n'est pas \unsure{constante|constant} sur les composantes \deleted{connexes} \inserted{irréd.} géom. de $X$, alors \deleted{\unsure{$g$}} on aurait $g(X(\overline K))=\overline K-\text{ens fini}$, donc \deleted{on aurait pour $v\in M_{\overline K}$} comme $v$ \inserted{$\in M_{\overline K}$} prend une infinité de valeurs, la propriété au-dessus ne serait pas vérifiée. Ainsi $g$ est algébrique sur $K$, et de plus en tant qu'\unsure{él\textsuperscript{mt}|élément|dit} de \inserted{la $K$-}algèbre séparable des sections \deleted{séparables} de $\underline O_X^{\times}$ alg. sur $K$, cela doit être une $M_K$-\uline{unité}.

Dans le cas particulier où $X$ est géom. intègre, on voit donc que le noyau de $\underline R_X^{\times}\unsure{-|\to}\Phi_X^*$ est le faisceau \deleted{\unsure{qui}} constant $U_X$ défini par le groupe $U$ des $M_K$-unités de $K$. On a donc les suites exactes can.
% archive sheet number: 43
\scanpage{45}
\diagram{\begin{tikzcd}[ampersand replacement=\&,column sep=small]
0\arrow[r] \& \underline{\Phi}_X\arrow[r] \& \underline{\mathrm{QF}}_X\arrow[r,very thick] \& \underline{\mathrm{Div}}_X\arrow[r,dash] \& 0
\end{tikzcd}}
\diagram{\begin{tikzcd}[ampersand replacement=\&,column sep=small]
0\arrow[r] \& \underline R_X^*/\underline U_X\arrow[r] \& \underline{\mathrm{QF}}_X\arrow[r] \& \underline{\Phi}_X/\operatorname{Im}\underline O_X^*\arrow[r] \& 0
\end{tikzcd}}

Qu'on a avantage à envisager en diagramme \unsure{commun|comme|commutatif}.

\diagram{\footnotesize\begin{tikzcd}[ampersand replacement=\&,column sep=small,row sep=1.4em]
\& \& 0 \& \& \\
0\arrow[r,phantom,"{\simeq}" description] \& \underline{\Phi}_X/\underline{\operatorname{Im}}\underline O_X^*\arrow[r,phantom,"{\simeq}" description] \& \underline{\Phi}_X/\underline{\operatorname{Im}}\underline O_X^*\arrow[u]\arrow[r,phantom,"{\simeq}" description] \& 0 \& \\
0\arrow[r,dash] \& \underline{\Phi}_X\arrow[u,dashed]\arrow[r] \& \underline{\mathrm{QF}}_X\arrow[u]\arrow[r] \& \underline{\mathrm{Div}}_X\arrow[u,dash,dashed,"{\unsure{\sim|?|\wr}}" sloped] \& \\
0\arrow[r] \& \underline O_X^*/\underline U_X\arrow[u]\arrow[r] \& \underline R_X^*/\underline U_X\arrow[u]\arrow[r,dashed] \& \underline{\mathrm{Div}}_X\arrow[u,dashed,"{\sim}" sloped]\arrow[r,dashed] \& 0\\
\& 0\arrow[u] \& 0\arrow[u] \& 0\arrow[u,dash,dashed,"{\sim}" sloped] \&
\end{tikzcd}}

\unsure{,|.|{}} Passant aux sections globales, on trouve $\bigl\{$ pour $X$ géom intègre [ce qui implique que $\underline R_X^*/\underline U_X$ est constant, donc flasque $\bigr\}$, on obtient le diagramme commutatif à lignes et colonnes exactes
% archive sheet number: 44
\scanpage{46}
\[
\deleted{\varphi:X(K)\times M_K\longrightarrow\mathbb R}
\]
\deleted{est dite \uline{loc.} $M_K$-\uline{bornée}, \inserted{ou \uline{admissible}} si elle est \inserted{$M_K$-}bornée sur les parties $M_K$-bornées, \inserted{i.e. si pour tt tel $E$} Noter que \deleted{\unsure{une partie de}} \inserted{pour \unsure{tte|tt}} \deleted{partie finie $F$} \inserted{$P\in X(K)$,} \unsure{de} \deleted{$X(\overline K)\times X(\overline K)$}, $\{P\}\times M_K$ est $M_K$-borné, donc si $\varphi$ est loc. $M_K$-bornée, alors \deleted{\unsure{la}} \inserted{pour $P\in X(K)$,} $\varphi(P,\cdot):M_K\longrightarrow\mathbb R$ est un \deleted{$M_{\overline K}$} \deleted{$M_{\overline K}$-diviseur} $M_K$-diviseur. Donc une $\varphi$ admissible s'interprète aussi comme une application}
\[
\deleted{\widetilde{\varphi}:X(K)\longrightarrow\operatorname{Div}(M_K)}
\]
\deleted{qui a \inserted{tout au moins} la propriété suivante : Si $F$ est une partie $M_K$-bornée de $X(K)$ [i.e. telle que $F\times M_K$ soit $M_K$-borné], \uline{alors} l'ens des $\widetilde{\varphi}(P)$ $(P\in F)$ est \uline{borné} dans l'ens des diviseurs de $M_K$ [\uline{considéré comme ens ordonné}].}

\margin{\deleted{?}}
\deleted{\uline{Cette condition est-elle suffisante} ??}

\deleted{\deleted{\unsure{ou|on}} Une fonction $\varphi:X(K)\times M_K\longrightarrow\mathbb R$ qui est $M_K$-bornée est aussi dite \uline{\unsure{équivalente à zéro}} (elle est alors admissible)}
% archive sheet number: 45
\scanpage{47}
\diagram{\resizebox{\linewidth}{!}{\begin{tikzcd}[ampersand replacement=\&,column sep=small,row sep=1.5em]
\& 0 \& 0 \& 0 \& 0\\
0\arrow[r] \& \Phi(X)/(\mathrm G_m(X)/U)\arrow[u]\arrow[r] \& \Gamma(X,\underline{\Phi}_X/(O_X^*/U_X))\arrow[u]\arrow[r] \& \operatorname{Pic}(X)\arrow[u]\arrow[r,"{\deleted{\unsure{0\to H^1|\longrightarrow}}}" description] \& H^1(X,\underline{\Phi}_X)\arrow[u]\\
0\arrow[r,dash] \& \Phi(X)\arrow[u]\arrow[r] \& \mathrm{QF}(X)\arrow[u]\arrow[r] \& \operatorname{Div}(X)\arrow[u]\arrow[r] \& H^1(X,\underline{\Phi}_X)\arrow[u,phantom,"{=}" description,sloped]\\
0\arrow[r] \& \mathrm G_m(X)/U\arrow[u]\arrow[r] \& R(X)^*/U\arrow[u]\arrow[r] \& \operatorname{Div}_{\ell}(X)\arrow[u]\arrow[r] \& 0\arrow[u]\\
\& 0\arrow[u] \& 0\arrow[u] \& 0\arrow[u] \&
\end{tikzcd}}}

(utilisant que $\underset{\substack{\parallel\\\operatorname{Pic}(X)}}{H^1(X,\underline O_X^{\times})}\simeq H^1(X,\underline O_X^{\times}/U)$ grâce au fait que $U$ est flasque).

Je dis que \uline{si $X$ est quasi-projective}, alors dans le diagramme comm. précédent on peut remplacer $H^1(X,\underline{\Phi}_X)$ par 0 (deux fois), en d'autres termes que $\Gamma(X,\underline{\Phi}_X/(\underline O_X^{\times}/U_X))\to\operatorname{Pic}(X)$ est surjectif. Comme tt $\unsure{\xi|\zeta}\in\operatorname{Pic}(X)$ est la différence de deux éléments très amples, on est ramené à prouver que tt élément très ample est dans l'image, et
% archive sheet number: 46
\scanpage{48}
\deleted{Cela signifie aussi que $\widetilde{\varphi}(X(K))$ est une partie bornée de $\operatorname{Div}M_K$.}

\deleted{Si $K'$ est une ext. finie de $K$, alors on a $M_{K'}$, et $M_{K'}\longrightarrow M_K$. Toute fonction}
\[
\deleted{\varphi:X(\overline K)\times M_{K'}\longrightarrow\mathbb R}
\]
\deleted{définit une fonction}
\[
\deleted{\varphi:X(K)\times M_K\unsure{-|\longrightarrow}\mathbb R}
\]
% archive sheet number: 47
\scanpage{49}
par fonctorialité ceci nous ramène au cas où $X=\mathbb P^r_K$, $\unsure{\xi|\zeta}=\underline O_X(1)$.
% archive sheet number: 48
\scanpage{50}
\figure{\hfill\tikz[x=.52cm,y=.52cm]{\draw[thin] (0,0)--(-.12,1.6)--(.43,.03)--(.31,1.84)--(.82,.17) .. controls (.88,-.04) and (.8,-.12) .. (.77,-.13); \draw[thin] (.83,.67)--(.97,1.34) (1.04,.61)--(1.22,1.41);}}
\mspage{- 82 -}\hfill n\textsuperscript{o}357

\[
\deleted{f^\rho(x\sigma)=\chi(\sigma)^{-1}f^{\handwritten{\rho}}(x)}
\]
\deleted{\uline{quel que soit} $\sigma\in H$.}

\deleted{Soit $K$ compact dans $X$ ; alors $x\in K$, $x\xi\in S$ entraîne $x\xi\in S\cap KH$, d'où, puisque $S\cap KH$ est compact et que $H$ opère proprement, $\xi\in L$, où $L$ est une partie compacte de $H$ qui ne dépend que de $K$. Pour $x\in K$, on peut donc remplacer $\int_H$ par $\int_L$ dans (1) ; d'après le lemme 2 du Chap.III, \S\,5, n\textsuperscript{o}1, il s'ensuit que $f^{\handwritten{\rho}}$ est continue sur $K$. Le reste est évident.}

\deleted{Si une fonction $g$ satisfait à $g(x\sigma)=\chi(\sigma)^{-1}g(x)$ quels que soient $x\in X$, $\sigma\in H$, son support $S$ est invariant par $H$, et s'écrit donc $\pi^{-1}(\pi(S))$. On désignera par $\handwritten{\unsure{\mathscr K|\mathcal K}}^{\handwritten{\chi}}(X,H)$ l'ensemble des fonctions continues $g$ sur $X$, à valeurs dans $\underline{\underline{\mathrm R}}$, qui satisfont à $g(x\sigma)=\chi(\sigma)^{-1}g(x)$ quels que soient $x\in X$ et $\sigma\in H$, et dont le support est de la forme $KH$ avec $K$ compact ; il est clair que c'est un espace de Riesz ; pour $\chi=1$, $h\to h\circ\pi$ est un isomorphisme de $\handwritten{\mathscr K}(X/G)$ sur $\handwritten{\mathscr K}^{1}(X,G)$. La prop.1 montre que $f\to f^\rho$ détermine une application de $\handwritten{\mathscr K}(X)$ dans $\handwritten{\mathscr K}^{\handwritten{\chi}}(X,G)$, qui est manifestement \uline{croissante}. On désigne par $\handwritten{\mathscr K}_+^{\handwritten{\chi}}(X,H)$ l'ensemble des fonctions $\geq 0$ appartenant à $\handwritten{\mathscr K}^{\handwritten{\chi}}(X,H)$.}

\deleted{PROPOSITION 2.- \uline{Soit} $K$ \uline{une partie compacte de} $X$. \uline{Alors il existe une fonction} $h\in\handwritten{\mathscr K}_+(X)$ \uline{telle que}, \uline{pour toute fonction} $g\in\handwritten{\mathscr K}^{\handwritten{\chi}}(X,H)$ \uline{de support contenu dans} $KH$, \uline{on ait} $g=(hg)^{\handwritten{\rho}}$.}

\deleted{Soit $f\in\handwritten{\mathscr K}_+(X)$ telle que $f>0$ sur $K$ ; posons :}
\[
\deleted{f'(x)=\int_H f(x\xi)\chi(\xi)^{-1}\mathrm d\handwritten{\rho}(\xi)\ ;}
\]
\deleted{comme $\chi^{-1}\handwritten{\rho}$ est une mesure de Haar sur $H$, $f'$ est continue, d'après le cor. de la prop.1, et de la forme $f^b\circ\pi$, avec $f^b\in\handwritten{\mathscr K}(X/H)$. Soit $U$ l'ensemble ouvert des points où $f^b>0$ ; comme $f'>0$ sur $K$, on a $U\supset\pi(K)$ ; $f/f'$ est continue sur $\pi^{-1}(U)$. Soit $U'$ une partie ouverte de $X/H$, telle que $\pi(K)\subset U'$ et $\overline U'\subset U$ ; soit $F$ continue $\geq 0$ sur $X/H$, égale à 1 sur $\pi(K)$ et}
% archive sheet number: 49
